% !TEX program = xelatex
\documentclass{resume}
\usepackage[UTF8,fontset=windows]{ctex}

\begin{document}
\pagenumbering{gobble}

\name{周景潇}

\basicInfo{
  \email{2241470466@qq.com} \textperiodcentered\
  \phone{13626566112} \textperiodcentered\
  \github[github.com/juice094]{https://github.com/juice094}}

\section{\faGraduationCap\ 教育背景}
\datedsubsection{\textbf{甘肃农业大学}，信息科学技术学院}{2023.09 -- 2027.06}
数据科学与大数据技术 · 工学学士（预计 2027.06）

均分 83 / 核心课 86 · 专业前 10\% · 核心课程：数据结构(90)、大数据挖掘(97)、数据存储(96)、数据库(86)、操作系统(87)

\section{\faCogs\ 技能}
\begin{itemize}[parsep=0.5ex]
  \item \textbf{语言}：Rust（主力，3年，44+ workspace crate 经验）、Python（熟练）、Java/Go/TypeScript（熟悉）
  \item \textbf{AI / Agent}：ReAct + Plan 双模式 Agent、MCP 协议四传输全实现（stdio/SSE/HTTP/WebSocket）、RAG（BM25 + 向量混合检索 + RRF 融合）、Multi-Agent 调度、Candle GGUF 本地推理
  \item \textbf{后端 / 系统}：tokio 异步、Axum REST API、WebSocket、SPMC 事件总线、TLS 1.3（rustls）、Protobuf（prost）、NAT 穿透（STUN/UPnP/Relay）
  \item \textbf{数据库 / 存储}：SQLite（WAL / FTS5 / 自定义 UDF）、PostgreSQL、Redis、Tantivy BM25、sled
  \item \textbf{工程实践}：Cargo Workspace 架构设计、Contract-First、零 unwrap 工程基线、CI/CD（GitHub Actions）、TDD
\end{itemize}

\section{\faUsers\ 项目经历}
\datedsubsection{\textbf{Clarity} -- Rust 原生本地优先 AI Agent 运行时}{2026.04 -- 至今}
\role{Rust · 22 workspace crates · 1,889 tests · $\sim$150K LOC}{独立项目}
单二进制编排 LLM、MCP 工具、记忆系统与多 Agent 协作，六前端共享同一 Agent 内核。零 Python/Node.js/Ollama 外部依赖。
\begin{itemize}
  \item 设计 Contract-First 22-crate 分层架构：contract（零内部依赖）为契约根节点，TUI / 桌面 GUI / Web IDE / CLI / 系统托盘 / 移动端 FFI 六入口共享 Agent 内核
  \item 实现 ReAct + Plan 双模式 Agent 循环与四层审批链（Interactive/Smart/Plan/Yolo）：Plan Mode 将步骤遗漏率从 $\sim$40\% 降至接近零
  \item 构建混合记忆系统：SQLite FTS5 + BM25 + 向量搜索 + 四级压缩归档，RRF 融合算法，零外部向量数据库依赖
  \item 完整实现 MCP 协议四传输（stdio/SSE/HTTP/WebSocket），集成 6+ LLM Provider（OpenAI/Anthropic/DeepSeek/Kimi/Ollama/Candle GGUF），Provider 网格负载均衡与熔断
  \item 跨前端 SPMC 事件总线（WireMessage 统一协议），13 种 RenderLine 变体，各前端零互相 import
  \item 工程基线：1,889 测试全通过（1,554 lib + 275 bin + 34 doc + 26 集成），Clippy 零 warning，生产代码零 unwrap/expect/panic，12-job CI 全绿
\end{itemize}

\datedsubsection{\textbf{Devbase} -- 本地开发者工作空间世界模型编译器}{2026.05 -- 至今}
\role{Rust · 12 workspace crates + 主 crate · 71 MCP 工具 · 616+ tests}{独立项目}
将本地 Git 仓库、PARA 笔记、Skill 脚本与 YAML 工作流编译为 AI 可推理的结构化上下文。
\begin{itemize}
  \item 设计三层编译架构（感知 → 知识 → 策略）：tree-sitter 多语言解析 + Tantivy BM25 + SQLite 图存储
  \item 以 MCP 协议暴露 71 个 stdio 工具（5 stable + 58 beta + 8 experimental），统一 McpTool trait 与幂等路由
  \item 零云端依赖实现语义检索：SQLite BLOB + 自定义 cosine\_similarity UDF，配合 Candle/Ollama 本地 embedding
  \item 建立 G1-G7 / RF-1-RF-7 七条架构红线，CI 强制零生产 panic，schema 迁移自动备份
\end{itemize}

\datedsubsection{\textbf{Syncthing-Rust} -- P2P 文件同步守护进程}{2025.11 -- 至今}
\role{Rust · 13 workspace crates · 392 tests · $\sim$58K LOC · $\sim$13MB 单二进制}{独立项目}
Go Syncthing BEP 协议 wire 级兼容的 Rust 重新实现。生产部署于 Windows 11 ↔ Ubuntu 24.04 双节点。
\begin{itemize}
  \item 完整实现 BEP over TLS 协议栈：prost + rustls + ed25519-dalek，跨语言 wire\_compat 集成测试验证与 Go Syncthing 互操作
  \item 多路径 NAT 穿透：LAN UDP 广播 + Global HTTPS mTLS + STUN + UPnP + Relay v1，ParallelDialer 并行拨号
  \item 高可靠同步链路：SHA-256 块级扫描（rayon 并行）+ 自适应并发拉取（RTT 反馈）+ 三路文本合并 + Simple/Staggered 版本归档
  \item 零运行时依赖（无 OpenSSL/GC/Python），单静态二进制 $\sim$13MB，Clipy 零 warning
\end{itemize}

\section{\faFlask\ 研究经历}
\datedsubsection{\textbf{RAG 系统中输出结构的实证研究}}{2026.03 -- 至今}
独立完成。\href{https://github.com/juice094/acr-select}{github.com/juice094/acr-select}
\begin{itemize}
  \item 设计 $\sim$40 实验条件、3,000+ 样本的多因子受控实验，跨 7 模型 × 4 架构验证框架普适性
  \item 提出格式-内容双维度观测框架，主动发现测量 artifact 并设计统一评分协议修正统计推断
  \item 研究成果直接指导 Clarity / Devbase 中 RAG 系统的设计与调优
  \item 论文撰写中（manuscript in progress）
\end{itemize}

\section{\faCertificate\ 认证}
\begin{itemize}[parsep=0.5ex]
  \item AI 应用工程师（中级）--- 工业和信息化部认证
  \item 日语专业四级（72 分）
  \item C1 驾驶证
\end{itemize}

\section{\faInfo\ 其他}
\begin{itemize}[parsep=0.5ex]
  \item 求职意向：AI Agent 开发 / 后端开发 / 智能体平台研发
  \item 意向城市：上海 · 北京 · 深圳 · 杭州 · 成都
  \item 可到岗时间：2026.07；可实习 3--6 个月，每周 5 天
\end{itemize}

\end{document}
